\documentclass{article}
\usepackage{geometry}
\usepackage{graphicx} % Required for inserting images
\usepackage{amsmath, amsthm, amssymb}
\usepackage{parskip}
\newgeometry{vmargin={15mm}, hmargin={24mm,34mm}}
\theoremstyle{definition} 
\newtheorem{definition}{Definition}

\newtheorem{theorem}{Theorem}[section]
\newtheorem{lemma}[theorem]{Lemma}
\newtheorem{corollary}{Corollary}[theorem]


\title{MATH210B Homework 1}
\date{February 2024}
\author{Boran Erol}

\begin{document}

\maketitle

\section{Problem 2}

Let $ F $ be the field of fractions of $ R $.

For the forward direction, note that if a matrix $ a $ has a zero divisor in the ring
$ M_{n}(R) $, then it has a zero divisor in $ M_{n}(F) $, which means $ det(a) = 0 $
by linear algebra.

Similarly, if $ det(a) = 0 $ in $ R $, it's $ 0 $ in $ F $ which means there's a zero
divisor $ b \in M_{n}(F) $ by linear algebra. Then, by clearing the denominators of $ b $,
we get a zero divisor in $ M_{n}(R) $.
 
\section{Problem 9}

We'll use the following lemma in the problem:

\begin{lemma}
    Let $ I $ be an ideal of a non-zero commutative ring $ R $.
    Then, $ (R/I)[t] $ is isomorphic to $ (R[t]) / \hat{I} $,
    where $ \hat{I} $ in the second ring is the ideal generated by $ I $.
\end{lemma}
\begin{proof}
    Consider $ \phi : R[t] \xrightarrow{} (R/I)[t] $ mapping
    every coefficient modulo $ I $. $ \phi $ maps $ 1 $ to $ 1 $.
    $ \phi $ is clearly a ring homomorphism since it's simply applying the canonical
    projection map to every component. It's also clearly surjective. Thus, it suffices to show
    that $ \ker(\phi) = \hat{I} $ to conclude the proof.

    Let $ f \in \hat{I} $. Then, $ f = gx $ for some $ g \in R[t] $ and $ x \in I $.
    Then, $ x \mid C(f) $, so $ f \in ker(\phi) $.

    Let $ f \in \ker(\phi) $. Then, every coefficient of $ f $ is in $ I $.
    Then, every monomial of $ f $ is in $ \hat{I} $, so $ f $ is in $ \hat{I} $ since $ \hat{I} $
    is closed under addition.

    We thus conclude the proof.
\end{proof}

Recall that every nonzero commutative ring contains a maximal ideal $ M $.
Let $ M $ be a maximal ideal in $ R $.
Let $ \hat{M} $ be the ideal generated by $ M $ in $ R[t] $.
By the lemma above, we have that $ R[t] / \hat{M} $ is isomorphic to $ (R/M)[t] $.

Then, prime ideals of $ R[t] $ containing $ \hat{M} $ are in bijection with prime ideals
of $ (R/M)[t] $.

Thus, it suffices to show that every polynomial ring over a field $ F $ contains infinitely many prime ideals.

Since $ F[t] $ is a Euclidean domain using the degree of the polynomial, we can simply reuse Euclid's proof.

Suppose not. Let $ f_{1}, \ldots, f_{n} $ be the set of generators of the prime ideals.
Consider $ g = f_{1} \cdot \cdots \cdot f_{n}  + 1 $.
Since $ f_{i} $ divides $ g $ for some $ i \in [n] $, we have that $ f_{i} \mid 1 $, which is a contradiction,
since prime elements can't be units.

\end{document}
